\documentclass[a4paper,12pt]{article}

\usepackage[utf8]{inputenc}   % si compilation pdfLaTeX
\usepackage[T1]{fontenc}
\usepackage{lmodern}

\usepackage[french]{babel}
\usepackage{amsmath,amssymb}
\usepackage{fancyvrb}

\fvset{obeytabs=true, tabsize=3}
\newtheorem{ex}{Exemple}
\newcommand\set[1]{\ensuremath{\{#1\}}}

\begin{document}

\title{Module \emph{gestion\_bk\_joueur.pl}} 
\date{}
\maketitle

\section{A quoi sert il ?}
Gère la BK d’un joueur en autorisant les ajouts, retraits et modifications d’informations.

\section{Où est il ?}
	Dans le dossier \emph{programs}.
	
\section{Comment l’utiliser ?}
On l’insère comme module dans un programme swi-prolog :
	
\begin{Verbatim}
 :- use_module(gestion_bk_joueur).
\end{Verbatim}


\section{Prédicats disponibles}

\subsection{cree\_bk}

\begin{Verbatim}
cree_bk(Dossier, Joueur, 
		[
			Couleurs_croissantes,
			Couleurs_decroissantes,
			Ordre_couleurs_respecte,
			Somme_valeurs_respectee,
			Strategie_possible,
			Choix_valeurs_limites,
			Choix_valeurs_dichotomiques
		],
		Seuil,
		infos(TempsMoyenGlobal, TempsMoyenJoueur, NbCoupsMoyenGlobal, 
			   NbCoupsMoyenJoueur)
		)
\end{Verbatim}

Écrit dans le répertoire \emph{Dossier} une première BK intitulée \emph{bk\_joueur\_jean.pl} (si la variable \emph{Joueur} est instanciée avec \emph{jean}) en fonction 
des informations données en arguments.

L’appel de ce prédicat pourrait donner la BK suivante :
\begin{Verbatim}

/* ========================================================================== */
/*                              BK_PLAYER jean                                */
/* ========================================================================== */
/* ==================== INFORMATION ON THE PLAYER =========================== */

name(jean).

/* =================  CARACTERISTICS OF jean PLAYING ======================= */

respects_colours(jean).
respects_sum(jean).
dichotomic_player(jean).
fast_player(jean).

/* ============  PERSONAL SKILLS WITH RESPECT TO GAMES DIFFICULTIES ========= */

relative_game_difficulty(
			somme(16), 
			colors_order(red < green < blue = red), 
			solution([c(blue,5), c(yellow,5), c(green,4), c(red,2)]), 
			easy, medium
			 					).

\end{Verbatim}

\subsection{lit\_bk}

\begin{Verbatim}
lit_bk(
		Dossier, 				(input, String)
		Joueur, 					(input, String)
		En_tete, 				(output, liste de Strings)
		Caracteristiques, 	(output, liste de Strings)
		Strategies,				(output, liste de Strings)
		Difficulties			(output, liste de Strings)
	)
\end{Verbatim}

Relit la BK intitulée \emph{bk\_joueur\_jean} (si la variable \emph{Joueur} est instanciée avec \emph{jean}) située dans le répertoire \emph{Dossier}
et retourne les informations dans les 4 variables :
\begin{enumerate}
	\item \emph{En\_tete}, contient l’en tête de la BK :
		\begin{Verbatim}
		[
			/* ===============================	*/,
			/*       BK_PLAYER jean 			   */,
			/* ===============================	*/,
			/* == INFORMATION ON THE PLAYER	== */,
		]
		\end{Verbatim}
	\item \emph{Caracteristiques}, contient la liste des caractéristiques du joueur : 
		\begin{Verbatim}
			[name(jean)]
		\end{Verbatim}
	\item \emph{Strategies}, contient la liste des stratégies du joueur :
		\begin{Verbatim}
			[
				respects_colours(jean), 
				respects_sum(jean),
				dichotomic_player(jean), 
				fast_player(jean)
			]
		\end{Verbatim}
		\item \emph{Difficulties}, contient la liste des difficultés relatives du joueur en fonction des caractéristiques des parties
		\begin{Verbatim}
			[
				relative_game_difficulty(
					somme(16), 
					colors_order(red < green < blue = red), 
					solution([c(blue,5), c(yellow,5), c(green,4), c(red,2)]), 
					easy, medium)
			]
		\end{Verbatim}
\end{enumerate}

\subsection{ecrit\_bk}

\begin{Verbatim}
ecrit_bk(
		Dossier, 				(input, String)
		Joueur, 					(input, String)
		En_tete, 				(input, liste de Strings)
		Caracteristiques, 	(input, liste de Strings)
		Strategies,				(input, liste de Strings)
		Difficulties			(input, liste de Strings)
	)
\end{Verbatim}

Écrit (ou remplace) la BK intitulée \emph{bk\_player\_jean} (si la variable \emph{Joueur} est instanciée avec \emph{jean}) dans le répertoire \emph{Dossier}
(\emph{En\_tete}, \emph{Caracteristiques}, \emph{Strategies} et \emph{Difficulties} sont décrits en section précédente). 

\subsection{modifie\_bk}

\begin{Verbatim}
modifie_bk(
				Dossier, 						(input, String)
				Joueur, 							(input, String)
				Modifs_caracteristiques, 	(input, liste d’atomes)
				Modifs_strategies				(input, liste d’atomes),
				Modifs_difficulties			(input, liste d’atomes)
			)
\end{Verbatim}

Modifie la BK intitulée \emph{bk\_player\_jean} (si la variable \emph{Joueur} est instanciée avec \emph{jean}) située dans le répertoire \emph{Dossier}
en fonction des modifications présentes dans \emph{Modif\_caracteristiques} et \emph{Modif\_strategies} (l’en tête actuel de la bk est conservé).
 \emph{Modifs\_caracteristiques}, \emph{Modifs\_strategies} et \emph{Modifs\_difficulties} sont des listes d’atomes parmi :
 \begin{enumerate}
 	\item \emph{a(new\_atom)}: ajoute l’atome \emph{new\_atom},
 	\item \emph{r(atom)} : retire l’atome \emph{atom},
 	\item \emph{s(atom, new\_atom)} : remplace l’atome \emph{atom} par \emph{new\_atom}.
 \end{enumerate}
 
 Par l’exemple :
 \begin{Verbatim}
modifie_bk(
	‘../bk’, 						
	‘jean’, 						
	[a(age(jean, 31)), s(level(jean, 10), level(jean, 9))], 	
	[r(respect_sum(jean))],
	[s(difficulty(5), difficulty(7))]
			)
\end{Verbatim}
modifie la bk \emph{bk\_player\_jean} située dans le dossier \emph{../bk} :
\begin{enumerate}
\item ajoute l’atome \emph{age(jean, 31)} aux caractéristiques de \emph{jean},
\item remplace l’atome \emph{level(jean,10)} par \emph{level(jean, 9)} aux caractéristiques de \emph{jean},
\item retire l’atome \emph{respect\_sum(jean)} aux stratégies de jeu de \emph{jean},
\item substitue l’atome \emph{difficulty(5)} par \emph{difficulty(7)} aux difficultés relatives de \emph{jean}
\end{enumerate}

Il est possible de ne modifier qu’une partie de la BK en utilisant les prédicats suivants :

\begin{Verbatim}
modify_bk_characteristics(
				Dossier, 						(input, String)
				Joueur, 							(input, String)
				Modifs_characteristics 		(input, liste d’atomes)
			)
\end{Verbatim}



\begin{Verbatim}
modify_bk_strategies(
				Dossier, 						(input, String)
				Joueur, 							(input, String)
				Modifs_strategies				(input, liste d’atomes)
			)
\end{Verbatim}

\begin{Verbatim}
modify_bk_difficulties(
				Dossier, 						(input, String)
				Joueur, 							(input, String)
				Modifs_difficulties 			(input, liste d’atomes)
			)
\end{Verbatim}


\subsection{display\_bk}

\begin{Verbatim}
display_bk(
				Dossier, 						(input, String)
				Joueur, 							(input, String)
			)
\end{Verbatim}

Affiche la BK intitulée \emph{bk\_player\_jean} (si la variable \emph{Joueur} est instanciée avec \emph{jean}) située dans le répertoire \emph{Dossier}.

\end{document}